\chapter{Lire les valeurs}

Quelques conventions valent pour toutes les vues ; une fois que vous les connaissez, vous lisez n'importe quel
tableau de bord.

\section{Ce que vous voyez est l'état exact}
Nous n'affichons aucune moyenne et aucun lissage. Une valeur moyennée est une valeur que la puce n'a jamais eue, et
elle cache justement ce qui nous intéresse : un arpège qui change de note à chaque image, une enveloppe qui
redémarre. Si une valeur scintille, c'est que la musique la fait scintiller.

\section{Notes et fréquences}
Les notes sont écrites à l'anglaise : C, D, E, F, G, A, B, avec \# pour les dièses. L'octave 4 contient le la à
440~Hz (A4), le do du milieu du piano est C4. Quand la fréquence tombe entre deux notes, nous affichons la plus
proche ; la fréquence en hertz vous donne l'écart. Un léger désaccord entre deux voix qui jouent la même note est
souvent voulu : c'est ce qui épaissit le son.

\section{Niveaux}
Les niveaux sont exprimés par rapport au maximum de la puce : 0~dB est le plein volume, $-6$~dB la moitié de
l'amplitude, $-\infty$ le silence. Les barres suivent la même échelle.

\section{Registres}
Les registres sont en hexadécimal, précédés de \$, comme dans les documentations d'époque : \$D418 sur le SID,
\$804 sur la PC Engine. Nous les montrons tels qu'ils ont été écrits par le driver, pas tels que nous les
interprétons ; si vous cherchez dans une documentation, vous retrouverez les mêmes valeurs.
